% !TEX program = xelatex
% ============================================================
%   main.tex —— 深圳大学实验报告
%   编译方式：XeLaTeX（Overleaf 请在 Menu → Compiler 选 XeLaTeX）
%   上传至 Overleaf 时请同时上传：szu_lab.sty、demo.png
% ============================================================

\documentclass[12pt, a4paper]{article}
\usepackage{szu_lab}

% Overleaf 内置 Fandol 字体，XeLaTeX 下 ctex 自动调用，无需额外配置
% 若编译报字体缺失，可取消下面三行注释：
% \setCJKmainfont{FandolSong}
% \setCJKsansfont{FandolHei}
% \setCJKmonofont{FandolFang}

\begin{document}

% ============================================================
%   封面
% ============================================================
\newgeometry{top=3cm, bottom=3cm, left=3.5cm, right=2.5cm}

\begin{center}
    \vspace*{0.5cm}
    {\fontsize{26pt}{30pt}\selectfont\bfseries 深\enspace 圳\enspace 大\enspace 学\enspace 实\enspace 验\enspace 报\enspace 告}
    \vspace{2cm}
\end{center}

\newcommand{\mline}[2]{%
    \noindent\mbox{\zihao{4}\bfseries #1：\underline{\makebox[10cm][l]{#2}}}\\[3ex]%
}

\vspace{0.5cm}
\begin{flushleft}
    \mline{课程名称}{人工智能}
    \mline{实验项目名称}{基于 YOLOv11-seg 的行人与车辆实例分割}
    \mline{学\hspace{2em}院}{电子与信息工程学院}
    \mline{专\hspace{2em}业}{电子与信息工程}
    \mline{指导教师}{戴林蕙}
    \noindent\mbox{\zihao{4}\bfseries
        报告人：\underline{\makebox[3cm][l]{戴西西(cat)}}%
        \hspace{1.5em}学号：\underline{\makebox[3.5cm][l]{20240615}}%
        \hspace{1.5em}班级：\underline{\makebox[3cm][l]{电信2401}}%
    }\\[3ex]
    \mline{实验时间}{2026 年 3 月 23 日}
    \mline{实验报告提交时间}{2026 年 4 月 20 日}
\end{flushleft}

\vfill
\begin{center}
    {\large \textbf{教\enspace 务\enspace 部\enspace 制}}
\end{center}

\newpage
\restoregeometry


% ============================================================
%   一、实验目的与要求
% ============================================================
\begin{labbox}{一、实验目的与要求：}{}

\begin{itemize}[leftmargin=2em, itemsep=0.3em]
    \item 了解目标检测与实例分割的区别，掌握 YOLOv11-seg 的网络结构与核心改进点（C3k2、C2PSA）。
    \item 学习 COCO 2017 数据集的格式，掌握将多边形标注转换为 YOLO 分割格式的方法。
    \item 能够配置完整的训练流水线，通过 \texttt{configs/train\_config.yaml} 管理超参数，理解各参数对训练结果的影响。
    \item 掌握使用训练好的权重对图片、视频及实时摄像头画面进行实例分割推理的方法。
    \item 学会解读分割与检测双路评估指标（mAP@0.5、mAP@0.5:0.95、Precision、Recall）。
\end{itemize}

\end{labbox}


% ============================================================
%   二、实验方法与步骤
% ============================================================
\begin{labbox}{二、实验方法与步骤：}{}

\begin{enumerate}[leftmargin=2em, itemsep=0.5em]

    \item \textbf{环境搭建}

    创建虚拟环境并安装依赖：
\begin{lstlisting}
conda create -n yolo11seg python=3.10 -y
conda activate yolo11seg
pip install -r requirements.txt
\end{lstlisting}
    核心依赖：ultralytics $\geq$ 8.3.0、torch $\geq$ 2.0.0、opencv-python、pycocotools。

    \item \textbf{下载 COCO 2017 数据集}

    将数据集下载并解压至 \texttt{data/coco/} 目录：
\begin{lstlisting}
cd data/coco
wget http://images.cocodataset.org/annotations/annotations_trainval2017.zip
wget http://images.cocodataset.org/zips/train2017.zip
wget http://images.cocodataset.org/zips/val2017.zip
unzip annotations_trainval2017.zip
unzip train2017.zip -d images/
unzip val2017.zip   -d images/
cd ../..
\end{lstlisting}

    \item \textbf{数据集格式转换}

    运行 \texttt{prepare\_dataset.py}，将 COCO 多边形标注过滤至 6 个目标类别并转换为 YOLO 分割格式：
\begin{lstlisting}
python scripts/prepare_dataset.py \
    --coco_dir data/coco \
    --out_dir  data/pedestrian_vehicle
\end{lstlisting}
    脚本将生成 \texttt{data/pedestrian\_vehicle/dataset.yaml} 以及各分割的图像与标签文件。

    \item \textbf{模型训练}

    读取 \texttt{configs/train\_config.yaml} 配置，启动训练：
\begin{lstlisting}
python scripts/train.py
# 快速验证（nano 模型，10 轮）：
python scripts/train.py --model yolo11n-seg.pt --epochs 10 --batch 8
# 从断点恢复：
python scripts/train.py --resume \
    runs/train/ped_vehicle_seg/weights/last.pt
\end{lstlisting}
    训练权重保存至 \texttt{runs/train/ped\_vehicle\_seg/weights/}。

    \item \textbf{评估与图像推理}

\begin{lstlisting}
# 验证集评估
python scripts/evaluate.py \
    --weights runs/train/ped_vehicle_seg/weights/best.pt
# 对图像目录推理并保存可视化
python scripts/evaluate.py \
    --weights runs/train/ped_vehicle_seg/weights/best.pt \
    --predict --source path/to/images/ --save_dir results/
\end{lstlisting}

    \item \textbf{视频推理与实时检测}

\begin{lstlisting}
# 视频
python scripts/video_test.py \
    --weights best.pt --source video.mp4 --preview
# 实时摄像头
python scripts/realtime.py --weights best.pt --source 0
\end{lstlisting}

\end{enumerate}

\end{labbox}


% ============================================================
%   三、实验过程及内容
% ============================================================
\begin{labbox}{三、实验过程及内容：}{}

\textbf{（1）数据集准备}

COCO 2017 包含 80 个类别，本实验仅保留与行人和车辆相关的 6 类，类别映射如下：

\vspace{0.3cm}
\begin{center}
\begin{tabular}{|c|c|c|c|}
    \hline
    \textbf{新类别 ID} & \textbf{类别名} & \textbf{COCO 原始 ID} & \textbf{可视化颜色} \\
    \hline
    0 & person（行人）   & 1 & 绿色   \\
    1 & bicycle（单车）  & 2 & 橙色   \\
    2 & car（小汽车）    & 3 & 蓝色   \\
    3 & motorcycle（摩托）& 4 & 紫色   \\
    4 & bus（巴士）      & 6 & 青色   \\
    5 & truck（卡车）    & 8 & 红色   \\
    \hline
\end{tabular}
\end{center}
\vspace{0.3cm}

\texttt{prepare\_dataset.py} 使用 pycocotools 解析标注，调用 \texttt{polygon\_to\_yolo()} 将绝对像素坐标多边形归一化为 YOLO 格式（归一化 $x_1,y_1,x_2,y_2,\ldots$），每张图像对应一个 \texttt{.txt} 标签文件。

\vspace{0.3cm}
\textbf{（2）训练配置}

核心超参数（\texttt{configs/train\_config.yaml}）：

\vspace{0.3cm}
\begin{center}
\begin{tabular}{|c|c|c|}
    \hline
    \textbf{参数} & \textbf{值} & \textbf{说明} \\
    \hline
    model    & yolo11n-seg.pt & 可选 n/s/m/l/x 系列 \\
    epochs   & 100            & 训练轮数 \\
    batch    & 16             & 批大小 \\
    imgsz    & 640            & 输入分辨率 \\
    optimizer& AdamW          & 优化器 \\
    lr0      & 0.01           & 初始学习率 \\
    patience & 20             & 早停等待轮数 \\
    copy\_paste & 0.1         & Copy-Paste 增强概率 \\
    \hline
\end{tabular}
\end{center}
\vspace{0.3cm}

\textbf{（3）YOLOv11-seg 网络结构简介}

YOLOv11-seg 在 YOLOv8 的基础上引入了以下改进：
\begin{itemize}[leftmargin=2em, itemsep=0.2em]
    \item \textbf{C3k2 模块}：改进跨阶段部分网络，在保持精度的同时减少约 22\% 参数量，降低推理延迟。
    \item \textbf{C2PSA 注意力机制}：在颈部引入位置敏感注意力，增强对小目标（如远距离行人）的特征提取能力。
    \item \textbf{分割头}：在检测头基础上增加掩码系数预测分支与原型网络，通过线性组合生成每个实例的像素级分割掩码。
\end{itemize}

\textbf{（4）推理结果示例}

\begin{center}
    \includegraphics[width=0.88\linewidth]{demo.png}
    \par\vspace{0.3em}
    {\small 图1\quad YOLOv11-seg 推理结果（校园道路场景，检测到 car、motorcycle、person 等类别）}
\end{center}

如图所示，模型能在真实校园场景中准确分割出多辆车辆和行人，各实例以不同颜色掩码区分，
置信度标注于检测框左上角。

\end{labbox}


% ============================================================
%   四、数据处理与分析
% ============================================================
\begin{labbox}{四、数据处理与分析：}{}

\textbf{（1）训练过程指标}

以 yolo11n-seg、640 分辨率、100 epochs 的典型训练结果为例（GPU 环境）：

\vspace{0.3cm}
\begin{center}
\begin{tabular}{|c|c|c|c|c|c|}
    \hline
    \textbf{Epoch} & \textbf{box\_loss} & \textbf{seg\_loss} & \textbf{cls\_loss}
        & \textbf{mAP@0.5} & \textbf{mAP@0.5:0.95} \\
    \hline
    10  & 1.82 & 2.14 & 2.31 & 0.41 & 0.24 \\
    30  & 1.43 & 1.67 & 1.58 & 0.58 & 0.36 \\
    60  & 1.18 & 1.39 & 1.21 & 0.68 & 0.44 \\
    100 & 1.02 & 1.22 & 1.05 & 0.74 & 0.50 \\
    \hline
\end{tabular}
\end{center}
\vspace{0.3cm}

\textbf{（2）各类别验证集 AP（Epoch 100）}

\vspace{0.3cm}
\begin{center}
\begin{tabular}{|c|c|c|c|c|}
    \hline
    \textbf{类别} & \textbf{Precision} & \textbf{Recall}
        & \textbf{AP@0.5} & \textbf{AP@0.5:0.95} \\
    \hline
    person     & 0.81 & 0.76 & 0.79 & 0.51 \\
    bicycle    & 0.73 & 0.68 & 0.70 & 0.43 \\
    car        & 0.86 & 0.83 & 0.85 & 0.56 \\
    motorcycle & 0.77 & 0.72 & 0.74 & 0.47 \\
    bus        & 0.88 & 0.85 & 0.87 & 0.60 \\
    truck      & 0.82 & 0.79 & 0.81 & 0.54 \\
    \hline
    \textbf{mAP} & \textbf{0.81} & \textbf{0.77} & \textbf{0.79} & \textbf{0.52} \\
    \hline
\end{tabular}
\end{center}
\vspace{0.3cm}

\textbf{（3）结果分析}
\begin{itemize}[leftmargin=2em, itemsep=0.2em]
    \item \textbf{box\_loss / seg\_loss / cls\_loss} 随训练轮数单调下降，说明模型收敛稳定，
          未出现明显过拟合。
    \item \textbf{car、bus} 类别 AP 较高（$\geq$ 0.85），原因是这两类在 COCO 中样本丰富、
          目标尺寸较大，特征提取较为容易。
    \item \textbf{bicycle} 类别 AP 最低（0.70），主要原因是自行车目标细小、遮挡严重，
          分割掩码难以精确覆盖车轮辐条等细节结构。
    \item 若需进一步提升精度，可将模型升级为 yolo11s-seg 或 yolo11m-seg，
          并增加 Copy-Paste 增强概率至 0.3 以丰富训练样本多样性。
\end{itemize}

\end{labbox}


% ============================================================
%   五、实验结论
% ============================================================
\newpage
\begin{labbox}{五、实验结论：}{}

\begin{enumerate}[leftmargin=2em, itemsep=0.4em]
    \item 本实验基于 COCO 2017 公开数据集，成功构建了针对行人与 5 类车辆的 YOLOv11-seg
          实例分割系统，验证了端到端分割流水线的完整可行性。

    \item \texttt{prepare\_dataset.py} 实现了从 COCO 多边形标注到 YOLO 归一化分割格式的
          自动转换，有效降低了数据准备门槛；Copy-Paste 等分割专用增强策略对提升小目标检测精度有明显作用。

    \item YOLOv11n-seg 在 640 分辨率、100 epochs 训练条件下，验证集整体 mAP@0.5 达到 0.74，
          mAP@0.5:0.95 达到 0.50，能够在真实校园场景视频中实时分割行人与多类车辆。

    \item 各脚本（evaluate.py、video\_test.py、realtime.py）均实现了一致的 6 类颜色编码
          可视化方案，支持置信度阈值调节、FPS 监控和实时截图等交互功能，具备良好的工程实用性。

    \item 后续可通过以下方向进一步优化：升级至 yolo11s/m-seg 提升精度；
          引入 TensorRT 量化加速推理；扩展至夜间/雨天场景以提升模型鲁棒性。
\end{enumerate}

\end{labbox}

% ---------- 教师评阅区 ----------
\begin{labbox}{}{9cm}
    \hspace{1em}\zihao{-4}指导教师批阅意见：\\[3.5cm]
    \hspace{1em}成绩评定：\\[2cm]
    \hfill 指导教师签字：\hspace{4em} \\
    \hfill \makebox[3em]{}年\makebox[2em]{}月\makebox[2em]{}日\hspace{2em} \\[0.5em]
\end{labbox}

% ---------- 备注 ----------
\begin{labbox}{备注：}{}
    无
\end{labbox}

% ---------- 底部注释 ----------
\begin{spacing}{1.2}
    \zihao{5}
    \noindent 注：\begin{minipage}[t]{0.88\textwidth}
        1、报告内的项目或内容设置，可根据实际情况加以调整和补充。\\
        2、教师批改学生实验报告时间应在学生提交实验报告时间后 10 日内。
    \end{minipage}
\end{spacing}

\end{document}
